\documentclass[12pt,letterpaper]{article}
\usepackage[letterpaper,top=0.8in,bottom=0.8in,left=1in,right=1in]{geometry}
\usepackage[T1]{fontenc}
\usepackage{mathptmx}
\usepackage[hidelinks]{hyperref}
\usepackage{titlesec}
\setlength{\parindent}{0pt}
\setlength{\parskip}{5pt}
\AtBeginDocument{\setlength{\baselineskip}{13.7pt}}
\titleformat{\section}{\normalfont\normalsize\bfseries}{}{0pt}{}
\titlespacing*{\section}{0pt}{8pt}{3pt}
\hypersetup{pdftitle={Trustworthy Collective Reasoning and Evaluation for AI Agents},pdfauthor={Xinya Du}}
\begin{document}
\begin{center}
{\fontsize{15}{17}\selectfont\bfseries Trustworthy Collective Reasoning and Evaluation for AI Agents\par}
\vspace{4pt}
{\small Xinya Du | Assistant Professor of Computer Science | The University of Texas at Dallas\par}
{\small xinya.du@utdallas.edu | ONR FY27 Young Investigator Program | Program Officer Predrag Neskovic\par}
\end{center}

\section*{1 Problem}
How can a team of AI agents determine which judgments to trust and correct shared mistakes when the agents' expertise is uncertain? This project will develop methods for peer evaluation and collective reasoning that connect estimates of reliability to independent checks of disputed claims.

Consider an assistant helping a naval analyst assess whether a replacement component is available. Several agents endorse an answer using an old inventory report. Another finds a newer record and disagrees. The group must determine whether the new record applies to the same component and location, whether it is reliable, and whether it changes the answer. Counting endorsements does not resolve these questions. Agents using the same report may still contribute different reasoning, so shared sources alone also do not establish redundant judgments.

The scientific problem is to separate useful corroboration from shared error, then determine what additional evidence would resolve the uncertainty. We will focus on factual claims and reasoning steps whose correctness can be checked, at least in part, through records, tools, or independent audits.


\section*{2 Importance}
LLM-based assistants increasingly use other models to evaluate answers and critique reasoning. Their value depends on whether a second opinion adds information. A capable model may be a poor judge of a particular task, favor its own answer, or accept an unsupported argument. Agreement can then produce confidence without a corresponding improvement in accuracy.

For analysts combining maintenance, readiness, or mission-status reports, this creates a practical need: explain why an assessment deserves trust and identify the claim that still needs checking. The same research question arises in scientific reasoning and automated peer evaluation. Progress would help establish when collective AI judgments support a decision and when the available information calls for further review.


\section*{3 Current approaches and limitations}
LLM judges offer a scalable way to evaluate open-ended answers, although position and self-preference biases remain concerns \cite{ref1}. Peer ranking combines judgments from several models \cite{ref2}, while multi-agent debate allows models to inspect and revise one another's reasoning \cite{ref3}. These approaches provide useful starting points. However, a higher aggregate score or a successful discussion does not, by itself, identify which evidence made the final answer more reliable.

Computational social choice provides principles for selecting participants and allocating influence. Procaccia and collaborators study fair panel selection \cite{ref4}, impartial peer review in which a participant's ratings cannot change their own selection probability \cite{ref5}, and conditions under which delegation improves the chance of a correct collective decision \cite{ref6}. Work on manipulation-resistant selection also shows that strategic behavior requires explicit treatment \cite{ref7}. These results motivate precise objectives and safeguards for our peer-evaluation protocol.

Recent robust lotteries address a closely related AI-evaluation problem: aggregating heterogeneous pairwise preferences under shifts in prompts or annotators \cite{ref8}. Our focus is the correctness of claims and reasoning, including cases where several evaluators inherit the same factual mistake. A method can aggregate preferences consistently while still lacking the evidence needed to verify a claim.

The proposed contribution is to connect inference of peer reliability with verification during deliberation. A checked disagreement will update both the current answer and the evidence about which peers are reliable on similar claims. We will study what this feedback makes identifiable, how it should change agents' influence, and whether it improves later decisions.


\section*{4 Proposed research}
Our hypothesis is that collective reasoning improves when agents earn influence through demonstrated task expertise, and consequential disagreements trigger checks that can supply information beyond the group's existing judgments. The project has three connected components. Verification provides feedback to peer evaluation, and verified corrections provide supervision for learning.

\textbf{Learning which judgments deserve influence.} Agents will first produce answers and evaluations without seeing one another's judgments. The protocol will hide candidate identity where possible and retain the claims, cited sources, and available timestamps. We will extend peer ranking with a model of task-specific judge reliability and shared errors. Source overlap and model family will serve as clues about dependence; performance on independently checked claims will test whether those clues are useful.

A small set of audited claims will help distinguish a reliable evaluator from one that merely agrees with others. We will examine how much can be learned as the amount and coverage of audit data change. The weighting rule will limit concentrated influence from evaluators whose errors overlap, while allowing agents with complementary strengths to contribute. Review assignments and weight updates will also be designed to reduce an agent's ability to improve the standing of its own answer through favorable outgoing ratings. We will test residual vulnerability to coordinated ratings instead of assuming that blinding removes strategic behavior.

Analysis will begin with a simplified model of claim correctness, heterogeneous judge accuracy, and a shared source of error. We will characterize when distinct explanations of reliability produce the same observed peer judgments and what audited evidence can distinguish them. These results will define the assumptions under which a confidence estimate is justified. They will also identify cases that require an external check or abstention.

\textbf{Resolving consequential disagreements.} The deliberation method will align checkable claims across independently generated solutions and locate disagreements that could change the final answer. It will then select a check suited to the disputed claim. For a document-based claim, that may mean retrieving a record with different provenance and checking its date and scope. For a reasoning step, it may mean testing a counterexample or using an executable verifier.

Each check will record what it supports, what it contradicts, and what remains unresolved. That record will guide whether the group retains its answer, revises it, or requests review. The method will assess the substance of a challenge separately from how confidently an agent presents it. In the logistics example, finding a newer record is useful only after checking that it refers to the relevant component and depot. A correct minority should be able to change the answer through valid evidence; an unsupported minority should not gain influence merely by disagreeing.

\textbf{Learning from verified corrections.} We will retain the disputed claim, the check performed, and the resulting correction as feedback. This feedback will update which peers to consult on similar tasks and train policies for checking, revision, and abstention. We will also fine-tune open models on verified reasoning corrections, using process-level feedback to identify the step that needs repair. Training rewards will depend on independent verification or labeled outcomes, rather than agreement alone. This component will test whether a successful correction improves future decisions beyond the instance that produced it.

\textbf{Research foundation and feasibility.} Our group developed PRD for peer ranking and discussion \cite{ref2}, IQA-EVAL for interactive evaluation with simulated users \cite{ref9}, and FG-PRM for fine-grained detection of reasoning errors and solution verification \cite{ref10}. These works provide experience with the main experimental components: collecting peer judgments, evaluating interactions, and producing feedback on individual reasoning steps. The proposed work brings them together around a new question about reliability under dependence. The PI's 2024 NSF CAREER award also provides evidence of an independently led research program.

\textbf{Validation.} Experiments will combine multi-hop question answering and mathematical reasoning with controlled cases that vary shared sources, error correlation, and coordinated ratings. We will compare against a strong single judge, majority voting, unmodified peer ranking, and ordinary debate under comparable inference and checking budgets. Ablations will isolate reliability modeling, verification, and learning. The main outcomes are final-answer correctness, calibration, recovery of correct minority answers, and harmful revisions of initially correct answers. Human-preference agreement will be reported separately from factual correctness. Held-out topics and model families will test transfer. Reference answers used for evaluation will remain separate from the limited audits available to the agents.


\section*{5 Limitations}
Peer judgments cannot establish absolute correctness when all agents share a misconception and no independent evidence is available. Our methods will need to recognize this limit and leave some claims unresolved. Retrieved records can also be incomplete or wrong, and executable checks cover only certain kinds of reasoning. We will measure sensitivity to unreliable verification instead of treating every check as ground truth.

Reliability estimates may transfer poorly to new topics or model families, especially when audits cover only a narrow domain. Controlled manipulation tests will establish performance against specified behaviors; they will not establish resistance to arbitrary collusion. Simulated personas may change wording without adding independent expertise, so any benefit must be demonstrated empirically. The initial work will therefore emphasize checkable factual and reasoning tasks. Operational naval deployment would require domain data, expert validation, and integration beyond this research effort.


\section*{6 Expected impact}
If successful, the project will provide a clearer account of when peer evaluation adds reliable information, algorithms that use verified disagreements to improve collective answers, and training methods that carry those corrections into future reasoning. The scientific contribution will include explicit limits on what can be inferred from agreement and evidence for when verification overcomes those limits.

For naval information analysis, the resulting capability would help an analyst trace an assessment to its supporting evidence, identify a disputed claim, and understand why the system changed its answer or requested review. We will release research code, evaluation protocols, and controlled test cases to support further work on reliable collaborative AI.

\clearpage
\begingroup
\small
\setlength{\parskip}{0pt}
\begin{thebibliography}{10}
\setlength{\itemsep}{5pt}
\bibitem{ref1} L. Zheng et al. (2023). \href{https://papers.nips.cc/paper_files/paper/2023/hash/91f18a1287b398d378ef22505bf41832-Abstract-Datasets_and_Benchmarks.html}{Judging LLM-as-a-Judge with MT-Bench and Chatbot Arena}. NeurIPS 36, Datasets and Benchmarks.
\bibitem{ref2} R. Li, T. Patel, and X. Du (2024). \href{https://openreview.net/forum?id=YVD1QqWRaj}{PRD: Peer Rank and Discussion Improve Large Language Model based Evaluations}. Transactions on Machine Learning Research.
\bibitem{ref3} Y. Du, S. Li, A. Torralba, J. B. Tenenbaum, and I. Mordatch (2024). \href{https://proceedings.mlr.press/v235/du24e.html}{Improving Factuality and Reasoning in Language Models through Multiagent Debate}. ICML, PMLR 235:11733-11763.
\bibitem{ref4} B. Flanigan, P. Gölz, A. Gupta, B. Hennig, and A. D. Procaccia (2021). \href{https://doi.org/10.1038/s41586-021-03788-6}{Fair Algorithms for Selecting Citizens' Assemblies}. Nature 596:548-552.
\bibitem{ref5} D. Kurokawa, O. Lev, J. Morgenstern, and A. D. Procaccia (2015). \href{https://procaccia.info/wp-content/uploads/2020/03/impartial.pdf}{Impartial Peer Review}. IJCAI, pp. 582-588.
\bibitem{ref6} A. J. Berinsky et al. (2025). \href{https://doi.org/10.1287/mnsc.2023.02470}{Tracking Truth with Liquid Democracy}. Management Science 71(9):7821-7839.
\bibitem{ref7} B. Flanigan, J. Liang, A. D. Procaccia, and S. Wang (2024). \href{https://procaccia.info/wp-content/uploads/2023/08/sortmanip.pdf}{Manipulation-Robust Selection of Citizens' Assemblies}. AAAI 38, pp. 9696-9703.
\bibitem{ref8} H. Khalaf, F. P. Calmon, D. Halpern, A. D. Procaccia, I. Shapira, and S. L. Wang (2026). \href{https://procaccia.info/wp-content/uploads/2026/01/maxlot.pdf}{Robust AI Evaluation through Maximal Lotteries}. ICML.
\bibitem{ref9} R. Li, R. Li, B. Wang, and X. Du (2024). \href{https://proceedings.nips.cc/paper_files/paper/2024/hash/c6a23b26eaaefd187973658f5001f4fe-Abstract-Conference.html}{IQA-EVAL: Automatic Evaluation of Human-Model Interactive Question Answering}. NeurIPS 37.
\bibitem{ref10} R. Li, Z. Luo, and X. Du (2025). \href{https://aclanthology.org/2025.findings-emnlp.228/}{FG-PRM: Fine-grained Hallucination Detection and Mitigation in Language Model Mathematical Reasoning}. Findings of EMNLP, pp. 4247-4278.
\end{thebibliography}
\endgroup
\end{document}
